Em Cunha, p. 337: \enquote{o recém-instituído colegiado de pajés passou à discussão de temas como a hierarquia de especialistas, o encaminhamento de pacientes e outras questões de procedimento. Entre os temas discutidos, o principal era a reivindicação de que o Estado apoiasse, e na prática financiasse, o exercício  da medicina tradicional. O raciocínio era transparente: se o conhecimento médico krahô era considerado importante por uma faculdade de medicina, então devia ser tratado do mesmo modo que a assistência médica ocidental (pública). Devia haver instalações adequadas e os pajés e seus auxiliares deviam ser pagos pelo Estado. A mera sugestão dessa proposta horrorizou a EPM [Escola Paulista de Medicina] e deixou constrangida a representante do Ministério da Saúde.}